\maketitle
\tableofcontents

% ============================================================
\chapter*{Dok.~314 -- Das Gitter im Hilbertraum}
\addcontentsline{toc}{chapter}{Dok.~314 -- Das Gitter im Hilbertraum}
% ============================================================

\section*{Worum es geht}

Die FFGFT rechnet geometrisch auf T$^4$ mit dem D4-Gitter und
quantenmechanisch auf $\mathcal{H} = L^2(\text{T}^4)\otimes
\mathbb{C}^3$ (Dok.~230/282). Beim Übergang stellt sich eine
naheliegende Frage: \emph{Werden die Gitter dabei flach?} Bleibt von
Kusszahl~24, Packungsdichte und Trialität etwas übrig, oder ist im
Hilbertraum jedes Gitter dasselbe?

Die Antwort lautet: Der Torus ist ohnehin flach -- die Krümmung war
nie der Unterschied --, und die Gitterstruktur geht nicht verloren,
sondern \textbf{wandert ins Spektrum und in die Faser}. Dieses
Dokument rechnet das durch (drei Skripte, Kap.~J) und findet:

\begin{enumerate}
  \item \textbf{Das Spektrum trennt, was die Krümmung nicht trennt
        [K]:} Z$^4$ beginnt bei $|k|^2 = 1$ mit Entartung~8, D4 bei
        $|k|^2 = 2$ mit Entartung~24; die ungeraden Normen fehlen
        bei D4 vollständig. Die erste Anregung trägt exakt
        $24 = 12+12$ -- FCC-Hälfte und Wicklungshälfte getrennt
        adressierbar.
  \item \textbf{Die Trialität ist eingebaut, nicht angeheftet [K]:}
        $|\text{Aut}(D4)| = 1152 = |W(F_4)|$, Quotient zu $W(D4)$
        ist $6 = |S_3|$; unter den 80 Ordnung-3-Elementen erzeugen
        16 exakt den T$^4/Z_3$-Orbifold mit $|\det(1-A)| = 9$
        Fixpunkten und wirken auf jedem Modenorbit als
        $\mathbb{C}^3$-Zirkulant.
  \item \textbf{Ein Negativbefund mit Satzcharakter [K]:} $5$ teilt
        $1152 = 2^7\cdot3^2$ nicht, also enthält $\text{Aut}(D4)$
        \emph{kein} Element der Ordnung~5. Die Fünffachdrehung, aus
        der $\theta = 2/9$ stammt (Dok.~293), ist mit dem Gitter
        prinzipiell unverträglich -- die Phase gehört einer
        anderen Schicht.
  \item \textbf{Die Verdopplung ist dennoch vorbereitet [K]:}
        $24 = 8\times3 = 4\times6$; vier disjunkte Sechserblöcke mit
        genau dem $C_3$-Inhalt, den $3\oplus3'$ verlangt.
\end{enumerate}

Kein Teil der A-Serie; Status: Arbeitsdokument. Bezüge: Dok.~230/282
($\mathcal{H} = L^2(\text{T}^4)\otimes\mathbb{C}^3$), 285
(Verdopplung $6 = 3\oplus3'$), 291 (Kanalanalyse), 293
($\theta = 2/9$ ikosaedrisch), 311 (Vier auf Drei, $24 = 12+12$).

% ============================================================
\section*{Leseführung: der rote Faden}
% ============================================================

\textbf{Die Frage.} Im Hilbertraum wird das Gitter zur Indexmenge
der Fourier-Moden. Damit ist es keine Kugelpackung mehr -- in diesem
Sinn wird jedes Gitter \enquote{flach}. Die Frage ist, ob damit auch
die Information verschwindet.

\textbf{Die Antwort in einer Zeile.} \emph{Geometrie $\to$ Spektrum
$+$ Faser.} Die Entartungen des Laplace-Operators sind die
Thetareihe des Gitters; die Kusszahl wird zur Entartung der ersten
Anregung; die Trialität wird zur Orbifold- und Faserstruktur.

\textbf{Was daraus folgt.} Erstens sind zwei flache Tori am
Spektrum unterscheidbar (Kap.~B). Zweitens wird die
$\tilde T\cdot m = 1$-Struktur bei $R_4 \ne R_{123}$ im Spektrum
ablesbar, und die Niveauordnung ändert sich an genau zwei
Wurzelradien (Kap.~C). Drittens braucht die $Z_3$ nicht gesetzt zu
werden -- sie ist Gittersymmetrie (Kap.~E).

\textbf{Die Grenze, sauber gezogen.} Die Fünffachsymmetrie ist am
Gitter unmöglich (Kap.~F). Das trifft die Verdopplung nicht, weil
diese auf der \emph{Faser} lebt und die Schnittstelle beider Ebenen
genau $C_3$ ist -- besetzt durch die Trialität (Kap.~G). Und die
Container-Struktur gilt auf \emph{jeder} Schale, kann also keine
auszeichnen (Kap.~H) -- dieselbe Muster-gegen-Besetzung-Logik wie in
Dok.~293.

\textbf{Der Vorbehalt.} Gerechnet wird mit dem glatten Laplace,
also $D_f = 4$ exakt. Alle Zahlen sind damit \emph{nackt};
Kap.~I trennt, was das kostet: Zählungen und Verhältnisse bleiben
invariant, absolute Energien tragen $-1{,}33\,\%$, und ob die
Entartungen unter der fraktalen Störung aufspalten, ist nicht
gerechnet.

\medskip
\noindent\textbf{In einem Satz:} Im Hilbertraum wird das Gitter
flach, aber nicht folgenlos -- was die Geometrie trug, tragen
Spektrum und Faser, und die Arbeitsteilung zwischen beiden ist
exakt bestimmbar, solange man die glatte Näherung als solche
buchführt.

% ============================================================
\section*{Zeichenerklärung}
% ============================================================

Die Tabelle erklärt die Symbole dieses Dokuments. Größen, die nur an
einer Stelle auftreten, werden dort eingeführt.

\textbf{Notationshinweis.} In der Gittertheorie ist $\Lambda$ das
übliche Zeichen für ein Gitter (Conway/Sloane). Im FFGFT-Korpus ist
$\Lambda$ jedoch belegt -- als kosmologische Konstante, in
$\Lambda$CDM und insbesondere als Abschlussskala $\Lambda^{*} =
1/R_H^2$ (Dok.~312). Dieses Dokument verwendet deshalb
durchgehend $\Gamma$ für das Gitter und $\Gamma^{*}$ für das duale
Gitter. Ebenso wird $L^{*}$ hier \emph{nicht} verwendet, da es in
Dok.~312 die Zyklenlänge bezeichnet; die Kantenlänge des Torus
heißt hier $R_{123}$ bzw.\ $R_4$.

\begin{center}
\renewcommand{\arraystretch}{1.3}
{\small\setlength\tabcolsep{5pt}
\begin{longtable}{lp{9.6cm}}
\toprule
\textbf{Symbol} & \textbf{Bedeutung} \\
\midrule
\endfirsthead
\toprule
\textbf{Symbol} & \textbf{Bedeutung} \\
\midrule
\endhead
\bottomrule
\endfoot
\multicolumn{2}{l}{\textit{Gitter und Torus}} \\
$\Gamma$ & Gitter in $\mathbb{R}^4$; definiert den Torus
  T$^4 = \mathbb{R}^4/\Gamma$ (statt des in der Gitterliteratur
  üblichen $\Lambda$, siehe Notationshinweis) \\
$\Gamma^{*}$ & duales Gitter; trägt die zulässigen Wellenvektoren
  der Fourier-Moden \\
$D4$ & Wurzelgitter $D_4$: alle ganzzahligen $k$ mit
  $\sum k_i$ gerade; Kusszahl 24 \\
$\text{Z}^4$ & kubisches Gitter $\mathbb{Z}^4$; Kusszahl 8,
  exakt selbstdual \\
$D4^{*}$ & duales D4-Gitter; bis auf Skalierung wieder ein
  D4-Gitter \\
$R_{123}$, $R_4$ & Radien der drei Raumrichtungen bzw.\ der
  vierten Richtung \\
$r$ & Radienverhältnis $r = R_4/R_{123}$ (dimensionslos) \\
\midrule
\multicolumn{2}{l}{\textit{Spektrum}} \\
$k$ & Wellenvektor (Modenindex), $k \in \Gamma^{*}$ \\
$k_4$ & vierte Komponente von $k$; $k_4 = 0$ trennt die
  FCC-Hälfte von der Wicklungshälfte \\
$|k|^2$ & Norm; zugleich Eigenwert des Laplace-Operators auf dem
  isotropen Torus \\
$E(k)$ & Eigenwert bei Anisotropie,
  $E = k_1^2+k_2^2+k_3^2+(k_4/r)^2$ \\
$N(n)$ & Entartung der Schale $|k|^2 = n$; Koeffizient der
  Thetareihe \\
$\theta_{D4}$ & Thetareihe des D4-Gitters,
  $\theta_{D4} = (\theta_3^4+\theta_4^4)/2$ \\
$\theta(t)$ & Wärmekern-Thetafunktion, $\theta(t) = \sum_k
  e^{-t|k|^2}$ \\
$\zeta_M(s)$ & Epstein-Zetafunktion des Gitters $M$;
  $s = -1/2$ ist der Casimir-Punkt \\
\midrule
\multicolumn{2}{l}{\textit{Symmetrien}} \\
$\text{Aut}(\Gamma)$ & Automorphismengruppe des Gitters
  (längenerhaltende Selbstabbildungen) \\
$W(F_4)$, $W(D4)$ & Weylgruppen; $|W(F_4)| = 1152$,
  $|W(D4)| = 192$ \\
$A$ & Element der Ordnung 3 in $\text{Aut}(D4)$; halbzahlig, daher
  exakt über $2A$ gerechnet \\
$\det(1-A)$ & Determinante; $|\det(1-A)| = 9$ zählt die
  Orbifold-Fixpunkte \\
$Z_3$, $C_3$ & zyklische Gruppe der Ordnung 3 (Trialität bzw.\
  Faserwirkung) \\
$S_3$, $A_5$ & symmetrische Gruppe auf 3 Elementen; alternierende
  Gruppe auf 5 Elementen (ikosaedrisch) \\
$\omega$ & dritte Einheitswurzel, $\omega = e^{2\pi i/3}$;
  Phase $+120^\circ$ \\
$\chi_0,\chi_1,\chi_2$ & die drei $Z_3$-Charaktere; indizieren die
  Sektoren des Orbifolds \\
\midrule
\multicolumn{2}{l}{\textit{Hilbertraum und Faser}} \\
$\mathcal{H}$ & Zustandsraum,
  $\mathcal{H} = L^2(\text{T}^4)\otimes\mathbb{C}^3$
  (Dok.~230/282) \\
$\mathbb{C}^3$ & Faser über jedem Punkt; trägt die
  $Z_3$-Struktur \\
$3$, $3'$ & die beiden Tripletts der Verdopplung
  $6 = 3\oplus3'$ (Dok.~285) \\
$\varphi$ & Goldener Schnitt, $\varphi = (1+\sqrt5)/2$; tritt nur
  in der $A_5$-Schicht auf \\
$\theta = 2/9$ & Phasenanteil aus Dok.~293; keine Gitterinvariante
  (Kap.~F) \\
\midrule
\multicolumn{2}{l}{\textit{FFGFT-Bezüge}} \\
$\xi = 4/30000$ & geometrische Grundkonstante der FFGFT \\
$\tilde T\cdot m = 1$ & Zeit-Masse-Relation; im Spektrum als
  Aufspaltung bei $r \ne 1$ ablesbar \\
\midrule
\multicolumn{2}{l}{\textit{Statuslabel}} \\
{}[K] & Kernableitung: gerechnet und kontrolliert \\
{}[S] & Strukturaussage: konsistent, aber nicht abgeleitet \\
Satz & mathematisch bewiesen, nicht nur numerisch geprüft \\
Grenze & explizit gebuchte Reichweitenbeschränkung \\
\end{longtable}
}
\end{center}

% ============================================================
\chapter{Die Ausgangslage: Flachheit ist nicht der Unterschied}
% ============================================================

Der Torus T$^4 = \mathbb{R}^4/\Gamma$ ist \emph{immer} flach, egal
welches Gitter $\Gamma$ ihn definiert. D4-Torus und Z$^4$-Torus
haben beide Krümmung null. Der Unterschied liegt nicht in der
Krümmung, sondern in der Gitterstruktur: Kusszahl $24$ gegen $8$,
Packungsdichte $\pi^2/16$ gegen $\pi^2/32$.

Beim Übergang in $\mathcal{H} = L^2(\text{T}^4)$ wird das Gitter zur
\textbf{Indexmenge der Fourier-Moden}: Die Basis sind ebene Wellen
$e^{ik\cdot x}$, und die zulässigen $k$ liegen auf dem dualen Gitter
$\Gamma^{*}$. In diesem Sinn wird jedes Gitter im Hilbertraum
flach -- es ist keine Kugelpackung im Raum mehr, sondern ein
Punktgitter von Quantenzahlen.

\textbf{Die Geometrie geht dabei nicht verloren, sie wandert.} Der
Laplace-Operator hat Eigenwerte $|k|^2$, und die Entartung jedes
Eigenwerts ist die Anzahl der Gittervektoren dieser Länge -- die
Thetareihe des Gitters. Die 24 kürzesten D4-Vektoren werden zur
24-fachen Entartung der ersten Anregung, und die Zerlegung
$24 = 12+12$ aus Dok.~311 (FCC-Nachbarn plus Wicklungsnachbarn)
bleibt als $k_4$-Aufspaltung exakt lesbar.

Zwei Anker halten diese Übersetzung stabil:

\begin{enumerate}
  \item \textbf{D4 ist bis auf Skalierung selbstdual} -- $D4^{*}$
        ist wieder ein D4-Gitter. Die Fourier-Transformation führt
        nicht aus der Familie heraus. (Z$^4$ ist exakt selbstdual.)
  \item \textbf{Die Faser ändert am Gitter nichts.} In
        $\mathcal{H} = L^2(\text{T}^4)\otimes\mathbb{C}^3$ hängt die
        $\mathbb{C}^3$-Faser an jedem Punkt und trägt die
        $Z_3$-Struktur; das Modengitter bleibt $\Gamma^{*}$.
\end{enumerate}

\begin{equation}
  \boxed{\;\text{Geometrie}\;\longrightarrow\;
  \text{Spektrum}\;+\;\text{Faser}\;}
\end{equation}

% ============================================================
\chapter{Das Spektrum trennt, was die Krümmung nicht trennt [K]}
% ============================================================

Vollständige Schalenzählung bis $|k|^2 = 12$:

\begin{center}
\begin{tabular}{rrrc}
\toprule
Norm $|k|^2$ & Z$^4$ & D4 & D4-Split ($k_4{=}0$ / $k_4{\ne}0$) \\
\midrule
$1$ & $8$ & $0$ & --- \\
$2$ & $24$ & $24$ & $12$ / $12$ \\
$3$ & $32$ & $0$ & --- \\
$4$ & $24$ & $24$ & $6$ / $18$ \\
$5$ & $48$ & $0$ & --- \\
$6$ & $96$ & $96$ & $24$ / $72$ \\
$7$ & $64$ & $0$ & --- \\
$8$ & $24$ & $24$ & $12$ / $12$ \\
$9$ & $104$ & $0$ & --- \\
$10$ & $144$ & $144$ & $24$ / $120$ \\
$11$ & $96$ & $0$ & --- \\
$12$ & $96$ & $96$ & $8$ / $88$ \\
\bottomrule
\end{tabular}
\end{center}

Drei Ablesungen:

\textbf{(i)~Zwei flache Tori sind spektral unterscheidbar.} Z$^4$
beginnt bei Norm~1 mit Entartung~8, D4 bei Norm~2 mit Entartung~24.
Bei D4 fehlen die ungeraden Normen \emph{vollständig} -- die
Paritätsbedingung $\sum k_i$ gerade schließt sie aus. Was die
Krümmung nicht unterscheidet, unterscheidet das Spektrum sofort.

\textbf{(ii)~Unabhängige Kontrolle bestanden.} Die D4-Reihe
$24, 24, 96, 24, 144, 96 \dots$ ist die bekannte Thetareihe
$\theta_{D4} = (\theta_3^4 + \theta_4^4)/2$; die Skriptrechnung
reproduziert sie koeffizientenweise.

\textbf{(iii)~Die $12+12$-Zerlegung ist spektral adressierbar.} Die
erste Anregung trägt exakt zwölf Moden mit $k_4 = 0$ (die räumliche
FCC-Hälfte) und zwölf mit $k_4 \ne 0$ (die Wicklungshälfte). Bei
höheren Schalen ist die Aufspaltung asymmetrisch ($6/18$ bei Norm~4,
$8/88$ bei Norm~12): \textbf{Die $12/12$-Symmetrie ist eine
Eigenschaft der ersten Schale, kein Allgemeinsatz.} Das ist
ausdrücklich zu buchen -- Dok.~311 argumentiert mit $24 = 12+12$,
und diese Aussage gilt genau dort, wo sie gemacht wird.

% ============================================================
\chapter{Anisotropie: wo die Zeit-Masse-Relation ablesbar wird [K]}
% ============================================================

Wird die vierte Richtung mit anderem Radius kompaktifiziert
($r = R_4/R_{123}$, Eigenwert $E(k) = k_1^2+k_2^2+k_3^2+(k_4/r)^2$),
spaltet die erste Schale auf:

\begin{center}
\begin{tabular}{rll}
\toprule
$r$ & $E$ (Wicklung, 12 Moden) & $E$ (FCC, 12 Moden) \\
\midrule
$1{,}0$ & $2{,}000$ & $2{,}000$ (entartet: 24) \\
$1{,}1$ & $1{,}826$ & $2{,}000$ \\
$1{,}5$ & $1{,}444$ & $2{,}000$ \\
$2{,}0$ & $1{,}250$ & $2{,}000$ \\
\bottomrule
\end{tabular}
\end{center}

Die FCC-Moden bleiben bei $E = 2$ fest, die Wicklungsmoden folgen
exakt
\begin{equation}
  E_\text{Wicklung} = 1 + \frac{1}{r^2},
  \qquad
  \Delta = 1 - \frac{1}{r^2}.
\end{equation}
Die Aufspaltung ist ein direktes Maß für das Radienverhältnis --
\textbf{die Stelle, an der die $\tilde T\cdot m = 1$-Struktur im
Spektrum ablesbar wird}: Massenrichtung und Raumrichtungen tragen
verschiedene Anregungsenergien, sobald sie verschieden lang sind.

\section{Die Kreuzungskarte ist diskret und kurz}

Alle Niveaus haben die Form $E = a + b/r^2$ mit ganzzahligen
$a, b$. Im Bereich $1 < r \le 2{,}5$ gibt es unter allen Niveaus bis
Norm~8 \textbf{exakt zwei} Kreuzungsradien:

\begin{center}
\begin{tabular}{ll}
\toprule
$r = \sqrt{2}$ & reine Wicklung $(0,0,0,\pm2)$ kreuzt die FCC-12 \\
$r = \sqrt{3}$ & reine Wicklung kreuzt die Wicklungs-12 \\
\bottomrule
\end{tabular}
\end{center}

Die Niveaufolge ändert sich also nicht kontinuierlich-chaotisch,
sondern an zwei ausgezeichneten Wurzelradien. Jede Ordnungsaussage
über die niedrigsten Anregungen zerfällt in genau drei Regime:
$r < \sqrt2$, $\sqrt2 < r < \sqrt3$, $r > \sqrt3$. Jenseits von
$\sqrt3$ ist der \emph{Grundton} des Spektrums eine reine Wicklung
(bei $r = 2$: $E = 1{,}00$ mit 2 Moden vor $1{,}25$ mit 12 vor
$2{,}00$ mit 12).

% ============================================================
\chapter{Umorganisationen: drei Ebenen und eine Grenze}
% ============================================================

Flache Gitter lassen sich umorganisieren -- auf drei Ebenen, die
sehr Verschiedenes kosten.

\textbf{Basiswechsel (kostenlos).} Jede unimodulare Transformation
aus $GL(4,\mathbb{Z})$ beschreibt dasselbe Punktgitter mit anderen
Erzeugern. Das Spektrum ändert sich um nichts. Reine Buchhaltung.

\textbf{Gittersymmetrien (kostenlos, strukturtragend).}
Umorganisationen, die das Gitter auf sich abbilden -- hier sitzt der
zentrale Befund, Kap.~E.

\textbf{Umbau zu einem anderen Gitter (kostet Physik).} Diskret über
Klebevektoren entlang der Kette
\begin{equation}
  D4 \;\subset\; \text{Z}^4 \;\subset\; D4^{*}
  \qquad\text{(jeweils Index 2)},
\end{equation}
oder kontinuierlich durch Deformation im Modulraum
$GL(4,\mathbb{R})/(O(4)\times GL(4,\mathbb{Z}))$. Jede Stufe
verschiebt Entartungen; D4 ist darin ein ausgezeichneter Punkt
(Kusszahl 24, in 4D nachweislich nicht überbietbar).

\textbf{Die Grenze der Spektralaussage.} Ab Dimension~4 existieren
Paare flacher Tori, die \emph{isospektral, aber nicht isometrisch}
sind (Schiemann). Der Allgemeinsatz \enquote{Spektrum bestimmt das
Gitter} ist in 4D \emph{falsch}. Für D4 selbst unkritisch -- die
Kusszahl~24 zeichnet es eindeutig aus --, aber als Regel darf es
nicht gebucht werden.

% ============================================================
\chapter{Die Trialität ist eingebaut, nicht angeheftet [K]}
% ============================================================

\section{Am Gitter nachgerechnet}

\begin{center}
\begin{tabular}{lr}
\toprule
$|\text{Aut}(D4)|$ & $1152 = |W(F_4)|$ \\
$|\text{Aut}(\text{Z}^4)|$ & $384 = 2^4\cdot4!$ \\
$|\text{Aut}(D4)|/|W(D4)|$ & $1152/192 = 6 = |S_3|$ \\
\bottomrule
\end{tabular}
\end{center}

Beide Ordnungen wurden durch Brute-Force-Zählung aller
Basis-Tupel mit identischer Gram-Matrix bestimmt und treffen die
bekannten Werte. Der Quotient $6 = |S_3|$ weist die Trialität
\emph{direkt an der Gitter-Automorphismengruppe} nach, nicht nur am
Dynkin-Diagramm: D4 ist das einzige Gitter der $D_n$-Familie mit
drei gleichberechtigten Dynkin-Armen; bei $D5, D6, \dots$ schrumpft
die äußere Symmetrie auf $Z_2$.

\section{Die 80 Ordnung-3-Elemente in drei Klassen}

Exakt gerechnet (die Trialitätselemente sind \emph{halbzahlig} --
die Klebevektor-Mischung $(\frac12,\frac12,\frac12,\frac12)$ geht
ein; Rundung auf ganze Zahlen zerstört sie, deshalb wird mit $2A$
gerechnet):

\begin{center}
\begin{tabular}{lrl}
\toprule
Klasse & Anzahl & Charakter \\
\midrule
Spur $-2$, halbzahlig & $16$ & fixpunktfrei in $\mathbb{R}^4$;
  Eigenwerte $\{\omega,\bar\omega\}$ doppelt \\
Spur $1$, ganzzahlig & $32$ & Fixebene in $\mathbb{R}^4$,
  frei auf den 24 Minimalvektoren \\
Spur $1$, halbzahlig & $32$ & 6 Fixvektoren $+$ 6 Dreierorbits \\
\bottomrule
\end{tabular}
\end{center}

\section{Zwei Konsequenzen}

\textbf{(a)~Der Orbifold ist Gittersymmetrie, keine Zusatzannahme.}
Für jedes der 16 Elemente mit Spur $-2$ gilt exakt
\begin{equation}
  |\det(1-A)| = 9,
\end{equation}
der Quotient T$^4/\langle A\rangle$ ist also ein $Z_3$-Orbifold mit
\emph{genau 9 Fixpunkten} -- der Standardzahl $3^2$. Alle neun
liegen in Drittelkoordinaten. \textbf{Die T$^4/Z_3$-Struktur muss
nicht gesetzt werden; sie ist in D4 bereits vorhanden.}

\textbf{(b)~Auf jedem Modenorbit wirkt die $Z_3$ als Zirkulant.} Die
unitäre Wirkung auf $L^2(\text{T}^4)$ permutiert die drei ebenen
Wellen eines Dreierorbits zyklisch; die Permutationsmatrix hat
Eigenwerte $\{1,\omega,\omega^2\}$, Phasen $0^\circ$, $+120^\circ$,
$-120^\circ$ -- \emph{strukturell exakt die
Zirkulant-Diagonalisierung der $\mathbb{C}^3$-Faser} (Dok.~291). Die
erste Schale liefert acht Kopien: $24 = 8\times3
= 8\times(\chi_0\oplus\chi_1\oplus\chi_2)$.

% ============================================================
\chapter{Der Negativbefund: 5 teilt 1152 nicht [K]}
% ============================================================

Die Ordnungsverteilung in $\text{Aut}(D4)$, exakt bestimmt:

\begin{center}
\begin{tabular}{lrrrrrrr}
\toprule
Ordnung & $1$ & $2$ & $3$ & $4$ & $6$ & $8$ & $12$ \\
\midrule
Anzahl & $1$ & $139$ & $80$ & $228$ & $464$ & $144$ & $96$ \\
\bottomrule
\end{tabular}
\end{center}

Summe $1152$. \textbf{Ordnung 5: null Elemente} -- und das ist kein
Rechenergebnis, sondern ein Satz: $1152 = 2^7\cdot3^2$, und $5$
teilt das nicht (Lagrange).

\textbf{Konsequenz für $\theta = 2/9$.} Die Fünffachdrehung $R_5$,
aus der $p_0 = |\langle v_0|R_5 v_e\rangle|^2 = 2/9$ folgt
(Dok.~293), ist mit der D4-Gitterstruktur \emph{prinzipiell}
unverträglich. Der Test wurde zusätzlich direkt geführt: Alle
Invarianten der linearen Gitterwirkung an den neun Orbifold-Fixpunkten
($x\!\cdot\!x$, $k\!\cdot\!x$, $x\!\cdot\!y$, sämtlich mod 1 in exakter
Bruchrechnung) haben das Nennerspektrum $\{1,2,3,6\}$ --
\textbf{keine Neuntel}. Bemerkenswert dabei: $x\!\cdot\!y$ ist zwar
ein Neuntel-Vielfaches, aber die Zähler sind stets durch 3 teilbar;
die D4-Geometrie \emph{unterdrückt} Neuntel auf der ersten Stufe.

\textbf{Was das heißt -- und was nicht.} $2/9$ ist keine Invariante
der Gitterwirkung. Das entwertet Dok.~293 nicht, sondern grenzt es
ab: Die Phase gehört einer Schicht \emph{über} dem Gitter (der
$A_5$-Struktur), das Gitter trägt die $Z_3$. Die Arbeitsteilung ist
damit exakt benannt statt vermutet.

% ============================================================
\chapter{Die Verdopplung ist dennoch vorbereitet [K]}
% ============================================================

Widerspricht der Negativbefund der Verdopplung $6 = 3\oplus3'$ aus
Dok.~285? Nein -- in drei präzisierbaren Stufen.

\textbf{(1)~Das Verbot trifft die falsche Ebene.} $5\nmid1152$
verbietet Ordnung-5-\emph{Gitterautomorphismen}. Die Verdopplung
lebt auf der \textbf{Faser} ($\mathbb{C}^3\to\mathbb{C}^6$); $A_5$
wird nirgends als Raumsymmetrie gefordert. Die Kompatibilitäts\-bedingung
ist allein, dass die Schnittstelle beider Ebenen existiert -- und
die ist exakt bestimmbar: Elementordnungen in $A_5$ sind
$\{1,2,3,5\}$, in $\text{Aut}(D4)$ sind sie
$\{1,2,3,4,6,8,12\}$; der Schnitt ist $\{1,2,3\}$, kanonisch
also $C_3$. \textbf{Genau die stellt die Trialität bereit.} Die
Schnittstelle ist nicht nur erlaubt, sie ist die einzig mögliche --
und sie ist besetzt.

\textbf{(2)~Das Gitter bereitet die Verdopplung vor.} Die zentrale
Involution $-1 \in \text{Aut}(D4)$ kommutiert mit der
Trialitäts-$Z_3$ und paart die acht Dreierorbits der ersten Schale
in \textbf{vier disjunkte Paare}. Kein Orbit ist selbstgepaart, und
das beweisbar: $-k = A^j k$ erforderte den Eigenwert $-1$ für ein
Element der Ordnung~3 -- unmöglich. Also
\begin{equation}
  24 \;=\; 8\times3 \;=\; 4\times6
  \qquad\text{(Orbit $\oplus$ Spiegelorbit).}
\end{equation}
Jeder Sechserblock trägt unter $C_3$ den Charakterinhalt
$2\times(\chi_0\oplus\chi_1\oplus\chi_2)$ -- exakt die Restriktion
von $3\oplus3'$. Die erste Anregungsschale liefert vier fertige
Container.

\textbf{(3)~Die Grenze, sauber gebucht.} Die \emph{Etikettierung} --
welcher Faktor $3$ und welcher $3'$ ist -- ist gitterblind. Die
$A_5$-Charaktere der beiden Tripletts unterscheiden sich nur auf den
5er-Klassen ($\varphi$ gegen $1-\varphi$); auf der 3er-Klasse sind
beide null, die $C_3$-Restriktionen identisch. Die Unterscheidung
hängt exakt an der Struktur, die im Gitter nicht existiert --
vollständig konsistent mit Kap.~F.

\textbf{Bilanz:} Das Gitter liefert Container und $C_3$-Inhalt, die
$A_5$-Schicht liefert die $3/3'$-Etiketten, $\varphi$ und die
Fünffach -- und damit $2/9$.

% ============================================================
\chapter{Ankerlose Ergebnisse [K]}
% ============================================================

Alle folgenden Größen sind dimensionslos -- Verhältnisse, Zählungen,
Kongruenzen. Kein Kammerton, keine SI-Skala geht ein.

\section{Zwei Kongruenzsätze für die D4-Thetareihe}

Alle vollständigen Schalen bis Norm~120 geprüft:
\begin{equation}
  N(n) \equiv 0 \pmod 3
  \qquad\text{und}\qquad
  N(n) \equiv 0 \pmod 6
  \qquad\text{ausnahmslos.}
\end{equation}
Beides ist Satz, nicht Empirie: Die Trialitäts-$Z_3$ ist auf
$\Gamma\setminus\{0\}$ fixpunktfrei (Eigenwerte $\omega,\bar\omega$
-- kein Eigenwert 1), also zerfällt \emph{jede} Schale in
Dreierorbits; zusammen mit der $-1$-Paarung ohne Selbstpaarung in
Sechserblöcke.

\section{Charakteraufgelöstes T$^4/Z_3$-Spektrum}

Folge daraus: exakte Gleichverteilung auf jeder Schale.

\begin{center}
\begin{tabular}{rrrrr}
\toprule
Norm & $N(D4)$ & $\chi_0$ & $\chi_1$ & $\chi_2$ \\
\midrule
$2$ & $24$ & $8$ & $8$ & $8$ \\
$4$ & $24$ & $8$ & $8$ & $8$ \\
$6$ & $96$ & $32$ & $32$ & $32$ \\
$10$ & $144$ & $48$ & $48$ & $48$ \\
$14$ & $192$ & $64$ & $64$ & $64$ \\
$18$ & $312$ & $104$ & $104$ & $104$ \\
\bottomrule
\end{tabular}
\end{center}

Der invariante Sektor des Orbifolds ($\chi_0$) hat exakt $N(n)/3$
Zustände pro Schale; die $Z_3$-Projektion kostet auf keiner Schale
Asymmetrie.

\textbf{Ehrliche Konsequenz.} Gleichverteilung und Sechser-Container
sind \emph{keine Auszeichnung der ersten Schale}, sondern Satz für
alle. \textbf{Die Struktur ist überall -- und kann darum nichts
selektieren.} Das ist exakt die Muster-gegen-Besetzung-Logik aus
Dok.~293: Das Muster ist skaleninvariant präsent; welche Schale
physikalisch besetzt ist, entscheidet etwas anderes.

\section{Casimir: das dichteste Gitter verliert}

$\zeta$-regularisierte Vakuumenergie eines skalaren Feldes,
$E = \pi\,\zeta_M(-1/2)$, beide Tori auf Kovolumen~1 normiert
(Epstein-Zeta über die Riemann-Split-Darstellung, dreifach
verifiziert -- gegen Direktsumme bei $s = 4, 5$ und gegen die
geschlossene Form $Z_{\text{Z}^4}(s) = 8(1-4^{1-s})\zeta(s)\zeta(s-1)$,
Übereinstimmung auf 25 Stellen):

\begin{center}
\begin{tabular}{lrr}
\toprule
 & $\zeta_M(-1/2)$ & $E = \pi\zeta_M(-1/2)$ \\
\midrule
Z$^4$-Torus & $-0{,}2966893$ & $-0{,}9320768$ \\
D4-Torus & $-0{,}2764778$ & $-0{,}8685805$ \\
\midrule
Differenz & & $+0{,}0634963$ \\
\bottomrule
\end{tabular}
\end{center}

\textbf{Der Z$^4$-Torus liegt tiefer.} Der Grund ist die
Spektrallücke: Das dichteste Gitter D4 hat bei gleichem Volumen die
\emph{größte} erste Eigenfrequenz ($\sqrt2$ gegen $1$) und damit die
\emph{weniger} negative Nullpunktsenergie. Bemerkenswert ist die
Ordnungsumkehr: Für $s > 0$ minimiert D4 die Epstein-Zeta (die
bekannte Extremalität des dichtesten Gitters), bei $s = 0$ sind alle
Gitter gleich ($\zeta(0) = -1$ universell), und bei $s = -1/2$ ist
die Ordnung invertiert.

\textbf{Interpretationsgrenze, gebucht:} Das gilt für das skalare,
periodische Feld. Anderer Feldinhalt oder andere Randbedingungen
können das Vorzeichen der Bevorzugung ändern --
\enquote{die Natur wählt Z$^4$} folgt daraus \emph{nicht}.

\section{Wärmekern: die Beobachtbarkeitsgrenze beziffert}

Für beide Tori gilt $t^2\theta(t) \to 1$ (identischer Weyl-Term bei
gleichem Volumen); der Gitterunterschied ist exponentiell klein:

\begin{center}
\begin{tabular}{rr}
\toprule
$t$ & relative Differenz \\
\midrule
$0{,}5$ & $1{,}2\times10^{-2}$ \\
$0{,}2$ & $1{,}2\times10^{-6}$ \\
$0{,}1$ & $1{,}8\times10^{-13}$ \\
$0{,}05$ & $4{,}1\times10^{-27}$ \\
\bottomrule
\end{tabular}
\end{center}

\textbf{Jede Observable, die nur auf den führenden
Wärmekern-Koeffizienten beruht, ist gitterblind.} Der
D4/Z$^4$-Unterschied sitzt ausschließlich in der exponentiell
unterdrückten Thetareihen-Feinstruktur -- beziehungsweise direkt in
der Spektrallücke.

% ============================================================
\chapter{Was hier noch nackt ist [K]/offen}
% ============================================================

Dieses Dokument rechnet mit dem \emph{glatten} Laplace-Operator auf
$\text{T}^4 = \mathbb{R}^4/\Gamma$, also mit $D_f = 4$ exakt. Die
FFGFT trägt aber die fraktale Dimension $D_f = 4-\xi$ und den
Faktor $K_\text{frak} = 1-100\xi$ (A040). Sämtliche Zahlen hier sind
damit \textbf{nackt} -- in derselben Bedeutung, in der $E_0$ vor der
Klärung in Register~R72 nackt war. Die Konsequenz zerfällt in drei
sehr ungleiche Teile.

\section{Unberührt: alles Ankerlose}

Nach der Zuordnungsregel (Register~R70, Dok.~313) greift
$K_\text{frak}$ bei absoluten Größen und kürzt sich bei
Verhältnissen derselben Ebene. Damit bleibt unberührt:

\begin{itemize}
  \item alle \textbf{Zählungen} -- $24$, $8$, $N(n)$, $8\times3$,
        $4\times6$: ganzzahlige Kombinatorik, keine Länge;
  \item die \textbf{Kongruenzen} $N(n)\equiv0 \bmod 3$ und
        $\bmod\,6$: Gruppentheorie, kein metrischer Inhalt;
  \item die \textbf{Charakteraufteilung} $N/3$ je Sektor;
  \item alle \textbf{Verhältnisse} $E/E_\text{min}$ und die
        Kreuzungsradien $\sqrt2$, $\sqrt3$ -- Struktur durch
        Struktur, der Faktor kürzt sich heraus.
\end{itemize}

\noindent Das ist der eigentliche Grund, warum Kap.~H
\enquote{ankerlose Ergebnisse} heißt: Diese Aussagen sind
fraktal-invariant.

\section{Ein Faktor: alle absoluten Energien}

Mit einem Anker wird $E = \hbar H_0\,|k|$ (Grundmode $= \hbar H_0$,
siehe Kap.~B in Verbindung mit Dok.~312). Diese Werte sind nackt und
tragen den Faktor:

\begin{center}
\begin{tabular}{lrr}
\toprule
 & nackt & mit $K_\text{frak}$ \\
\midrule
erste D4-Schale ($|k| = \sqrt2$) & $3{,}230\times10^{-52}$~J &
  $3{,}187\times10^{-52}$~J \\
reine Wicklung ($|k| = 2$) & $4{,}567\times10^{-52}$~J &
  $4{,}507\times10^{-52}$~J \\
\bottomrule
\end{tabular}
\end{center}

\noindent Durchgehend $-1{,}33\,\%$. Dasselbe gilt für die
Casimir-Energien aus Kap.~H und für die Spektrallücke selbst.
Buchhalterisch erledigt, nicht strukturell interessant.

\section{Offen: der Operator, nicht die Umrechnung}

Der dritte Teil ist nicht durch Multiplikation zu erledigen. Bei
$D_f = 4-\xi$ ist der Operator \emph{kein glatter Laplace mehr} --
und exakte Symmetrie ist genau das, was die Entartungen trägt. Eine
fraktale Störung bricht sie:

\begin{center}
\begin{tabular}{lrr}
\toprule
Störungsstärke & $\Delta E/E$ & $\Delta E$ bei $E = \hbar H_0$ \\
\midrule
$O(\xi)$ & $1{,}3\times10^{-4}$ & $3{,}0\times10^{-56}$~J \\
$O(100\xi) = 1-K_\text{frak}$ & $1{,}3\times10^{-2}$ &
  $3{,}0\times10^{-54}$~J \\
\bottomrule
\end{tabular}
\end{center}

\noindent \textbf{Die 24-fache Entartung wäre dann keine Entartung
mehr, sondern ein Multiplett mit Feinstruktur.} In welche
Untermultipletts es zerfällt, entscheidet die \emph{Symmetrie der
Störung}:

\begin{center}
\begin{tabular}{ll}
\toprule
Störung respektiert & Aufspaltung der ersten Schale \\
\midrule
Trialität ($Z_3$) und $-1$ & $24$ bleibt zusammen \\
nur $Z_3$ & $8$ Dreierorbits bleiben zusammen \\
nur $-1$ & $4$ Sechserblöcke (Kap.~G) \\
nichts davon & bis zu $24$ Einzelniveaus \\
\bottomrule
\end{tabular}
\end{center}

\textbf{Gerechnet ist das nicht.} Es wäre die natürliche
Fortsetzung dieses Dokuments und keine große Rechnung --
Störungstheorie erster Ordnung auf einer endlichen Matrix, kein
Boltzmann-Apparat. Bis dahin gilt: Das Spektrum dieses Dokuments ist
die \emph{glatte Näherung}, und die Container-Struktur ist unter
Störung nur so stabil wie die Symmetrie, die sie trägt.

% ============================================================
\chapter{Statusübersicht}
% ============================================================

{\small
\begin{longtable}{lp{6.4cm}l}
\toprule
Aussage & Inhalt & Status \\
\midrule
\endfirsthead
\toprule
Aussage & Inhalt & Status \\
\midrule
\endhead
\bottomrule
\endfoot
Übersetzung & Geometrie $\to$ Spektrum $+$ Faser; Gitter wird
  Modenindex, Entartungen $=$ Thetareihe & [K] \\
Unterscheidbarkeit & Z$^4$ ab Norm 1 (8), D4 ab Norm 2 (24);
  ungerade Normen bei D4 leer & [K] \\
$12+12$ & erste Schale exakt FCC $+$ Wicklung; höhere Schalen
  asymmetrisch ($6/18$, $8/88$) & [K] \\
Anisotropie & $E_\text{Wick} = 1 + 1/r^2$; Kreuzungen nur bei
  $r = \sqrt2$ und $\sqrt3$; drei Regime & [K] \\
Trialität & $|\text{Aut}(D4)| = 1152$, Quotient $6 = |S_3|$;
  80 Ordnung-3-Elemente in drei Klassen & [K] \\
Orbifold & 16 Elemente mit $|\det(1-A)| = 9$: T$^4/Z_3$ ist
  Gittersymmetrie, keine Zusatzannahme & [K] \\
Zirkulant & $Z_3$ wirkt auf jedem Dreierorbit mit
  $\{1,\omega,\omega^2\}$ $=$ $\mathbb{C}^3$-Faserstruktur & [K] \\
$5 \nmid 1152$ & keine Ordnung-5-Automorphismen; $2/9$ ist keine
  Gitterinvariante (Nenner $\{1,2,3,6\}$) & [K] Satz \\
Verdopplung & Schnittstelle $A_5 \cap \text{Aut}(D4) = C_3$,
  besetzt; $24 = 4\times6$ als Container & [K] \\
Etikettierung & $3$ gegen $3'$ ist gitterblind (Unterschied nur auf
  5er-Klassen) & [K] Grenze \\
Kongruenzen & $N(n) \equiv 0 \bmod 3$ und $\bmod 6$ auf allen
  60 geprüften Schalen, mit Beweisgrund & [K] Satz \\
Selektion & Container und Gleichverteilung sind generisch --
  die Struktur kann keine Schale auszeichnen & [K] \\
Casimir & $E(\text{Z}^4) = -0{,}9321$ gegen
  $E(D4) = -0{,}8686$; Ordnungsumkehr gegenüber $s>0$
  & [K], Feldinhalt-abhängig \\
Wärmekern & Weyl-Term identisch, Differenz $4\times10^{-27}$ bei
  $t = 0{,}05$: führende Ordnung ist gitterblind & [K] \\
Spektralsatz & \enquote{Spektrum bestimmt Gitter} in 4D falsch
  (Schiemann); für D4 unkritisch, aber keine Regel & [K] Grenze \\
Fraktal-Status & Zählungen/Verhältnisse invariant; absolute Energien
  $-1{,}33\,\%$; Entartungsaufspaltung ungerechnet & [K]/offen \\
\end{longtable}
}

\medskip
\noindent\textbf{Bilanz:} Im Hilbertraum wird das Gitter flach, aber
nicht folgenlos. Die Kusszahl wird zur Entartung, die Trialität zum
Orbifold und zur Faser, das Radienverhältnis zur Niveauaufspaltung.
Was das Gitter \emph{nicht} trägt -- die Fünffachsymmetrie und damit
$\varphi$ und $2/9$ --, ist ebenso exakt bestimmt wie das, was es
trägt. Und die Container-Struktur ist generisch: Sie stellt bereit,
sie selektiert nicht.

% ============================================================
\chapter{Prüfskripte}
% ============================================================

Drei Skripte, alle deterministisch, alle Kontrollen bestanden:

\begin{itemize}
  \item \texttt{d4\_skript\_1\_spektrum\_deformation.py} --
        Thetareihen D4/Z$^4$, $k_4$-Split, Deformation $r$, exakte
        Kreuzungsradien (nur Standardbibliothek).
  \item \texttt{d4\_skript\_2\_trialitaet\_orbifold\_phasen.py} --
        $\text{Aut}(D4) = 1152$, Ordnung-3-Klassifikation (mit $2A$
        exakt gerechnet), $|\det(1-A)| = 9$, Zirkulant, Phasentest
        $2/9$, Ordnungsverteilung, $24 = 4\times6$ (numpy).
  \item \texttt{d4\_skript\_3\_schalen\_casimir.py} --
        Schalensätze, Epstein-Zeta bei $s = -1/2$, Wärmekern
        (numpy, mpmath).
\end{itemize}

Kontrollen gegen bekannte Werte: $\theta_{D4}$,
$|W(F_4)| = 1152$, $2^4\cdot4! = 384$, $|W(D4)| = 192$,
Orbifold-Fixpunktzahl $3^2 = 9$, Jacobi-Identität für
$Z_{\text{Z}^4}$.

\medskip
\noindent\textbf{Registrierung:} Aufnahme in Dok.~190 erfolgt
separat, sofern die Abgrenzung zu Dok.~293 dort gebucht werden soll.

\end{document}
